\chapter{Gérer les utilisateurs}

\section{Employés de l'organisation}

La page \emph{Users / Employees} permet de rechercher, filtrer, exporter et
administrer uniquement les employés rattachés à l'organisation courante.

\manualscreenshotcompact{admin-employees.png}
  {Gestion des employés d'une organisation}
  {fig:admin-employees}

\subsection{Créer un opérateur ou un technicien}

L'administrateur ne choisit pas le mot de passe de l'employé. Il crée une
invitation ; le destinataire active ensuite son propre compte.

\begin{procedurebox}{Inviter un employé}
  \begin{enumerate}[leftmargin=*]
    \item Ouvrir \emph{Users}, puis l'onglet \emph{Employees}.
    \item Choisir \emph{Add employee}.
    \item Saisir le nom, l'adresse électronique et sélectionner le rôle
    \emph{Operator} ou \emph{Technician}.
    \item Vérifier le quota du plan puis confirmer l'invitation.
    \item Contrôler que le compte apparaît avec un statut d'invitation en
    attente.
  \end{enumerate}
\end{procedurebox}

Le lien d'activation d'un employé est valide 72 heures par défaut. Cette durée
est configurable par le Super Admin entre 12 et 336 heures.

\subsection{Activation par l'employé}

Le destinataire ouvre le lien reçu, vérifie son identité, définit son mot de
passe puis complète les informations facultatives proposées. Une fois
l'activation terminée, le compte devient actif dans l'organisation et le rôle
choisi détermine immédiatement son espace de travail.

\expected{Le mot de passe n'est jamais communiqué à l'administrateur. Le
serveur stocke uniquement sa représentation sécurisée.}

\subsection{Renvoyer, révoquer ou désactiver}

\begin{itemize}[leftmargin=*]
  \item \textbf{Renvoyer} renouvelle le message d'une invitation encore valide
  ou expirée, sous réserve du délai minimal entre deux envois.
  \item \textbf{Révoquer} invalide définitivement le lien d'activation en
  attente.
  \item \textbf{Désactiver} bloque un compte déjà activé sans effacer son
  historique métier.
  \item \textbf{Réactiver} restaure l'accès d'un compte si le plan et
  l'organisation le permettent.
\end{itemize}

\attention{Ne pas supprimer un employé uniquement pour libérer un quota sans
avoir vérifié ses interventions, rapports et traces d'audit. La désactivation
préserve mieux l'historique.}

\section{Clients liés à l'organisation}

Un client possède un compte personnel global et peut utiliser les stations de
plusieurs organisations. Il devient visible pour une organisation lorsqu'une
relation métier existe, par exemple une adhésion, une session ou un paiement.
Le simple paiement d'une recharge ne crée pas automatiquement une adhésion à
un plan.

\begin{table}[H]
  \centering
  \caption{Différence entre employé et client}
  \label{tab:employee-customer}
  \small
  \renewcommand{\arraystretch}{1.08}
  \begin{tabularx}{\textwidth}{|L{3.2cm}|Y|Y|}
    \hline
    \rowcolor{ctmint}
    \textbf{Aspect} & \textbf{Employé} & \textbf{Client} \\ \hline
    Appartenance & Une seule organisation. & Compte global, relations possibles
    avec plusieurs organisations. \\ \hline
    Création & Invitation par l'administrateur. & Inscription locale publique ou
    connexion Google, selon les paramètres. \\ \hline
    Rôle & Opérateur ou technicien. & Conducteur utilisant le service. \\ \hline
    Quota commercial & Compte dans la capacité employés. & Ne compte pas comme
    employé. \\ \hline
    Administration & Gestion complète dans le périmètre de l'organisation. &
    Consultation des relations et activités autorisées, sans prise de contrôle
    du compte global. \\ \hline
  \end{tabularx}
\end{table}

\section{Export et confidentialité}

L'export d'une liste ne doit contenir que les données nécessaires à l'objectif
annoncé. Il faut protéger le fichier généré, limiter sa diffusion et le
supprimer lorsqu'il n'est plus utile.

\securitynote{Les mots de passe, secrets OAuth, jetons d'activation, clés API et
secrets OCPP ne doivent jamais apparaître dans un export, un rapport ou une
capture d'écran transmise au support.}

\begin{procedurebox}{Contrôler un export avant diffusion}
  \begin{enumerate}[leftmargin=*]
    \item vérifier que les filtres limitent l'export à l'organisation courante ;
    \item choisir le format adapté : PDF pour la lecture, CSV ou JSON pour le
    traitement ;
    \item ouvrir le fichier et contrôler les colonnes réellement présentes ;
    \item transmettre le document uniquement aux destinataires autorisés ;
    \item supprimer les copies temporaires après utilisation.
  \end{enumerate}
\end{procedurebox}
